% !TeX root = ../main.tex

Die in Abbildung~\ref{fig:s-rateExec} gezeigten Erfolgs- und Fehlerraten zeigen den Einfluss einer Strategie.
Für Deadlocks steigt die Erfolgsrate von \texttt{Strategie0} mit 18\% zu \texttt{Strategie1} 52\% extrem.
Dies ist auch die einzige Stelle, an der die Erfolgsrate höher ist als die Fehlerrate.
Bei \texttt{Strategie2} ist die Erfolgsrate wieder auf 46\% gesunken und bleibt gleich für \texttt{Strategie3}.
Die Erfolgsrate für Serialization Failures schaut ähnlich aus.
Erst steigt sie von 25\% für \texttt{Strategie0} auf 57\% für \texttt{Strategie1} stark an und schwächt ab und steigt
nur noch leicht an auf 60\% für \texttt{Strategie2}.
\texttt{Strategie3} hat mit 54\% eine geringere Erfolgsrate als die anderen beiden Retry Strategien.
Ab der ersten Strategie bleibt die Erfolgsrate über der Fehlerrate.
Lock Timeouts haben die niedrigste Erfolgsrate die für \texttt{Strategie0} bei 5\%.
Diese steigt dann mit \texttt{Strategie1} auf 12\%, mit \texttt{Strategie2} auf 15\% und sinkt leicht mit
\texttt{Strategie3} auf 14\%.
Die Ausführungszeiten von Deadlocks beginnen bei 2039ms und 2344ms und steigen im Laufe der Strategien.
Bei nicht erfolgreichen Transaktionen hat \texttt{Strategie1} den höchsten Wert mit 12100ms, danach sinkt dieser
wieder.
Erfolgreiche Transaktionen haben die höchste Ausführungszeit auch bei \texttt{Strategie2}.
Die Ausführungsdauer erfolgreicher Transaktion bei Lock Timeouts verhält sich wie bei den anderen Parametern.
Sie steigt leicht linear an bei knapp über 1000ms.
Für nicht erfolgreiche Transaktionen steigt die Dauer stark an von \texttt{Strategie0} zu \texttt{Strategie1}.
Danach steigt sie leicht weiter an und erreicht den Höchstwert bei \texttt{Strategie2} mit 573ms und sinkt bei
\texttt{Strategie3} wieder leicht ab.
Serialization Failures haben mit 52ms-54ms die niedrigste Ausführungszeit in \texttt{Strategie0}.
Zu \texttt{Strategie1} steigen beide Zeiten gleich an und trennen sich danach.
Nicht erfolgreiche Transaktionen haben bei \texttt{Strategie2} 263ms und dies ist damit der höchste Wert.
Für erfolgreiche Transaktionen liegt der höchste Wert mit 182ms bei \texttt{Strategie2}

\input{figures/s-rateExec}

Bei der Retry Verteilung~\ref{tab:s-retryDistribution} gibt es keine Auffälligkeiten, aber es ist wichtig zu erwähnen,
da der \emph{No Retry} wieder nicht auftaucht, da dieser unter \texttt{Strategie0} fällt.
Bei den Quantilen~\ref{tab:s-quantiles} ist erwähnenswert, dass bei Strategie das $p_{99.9}$-Quantil eine doppelt so
hohe Ausführungszeit wie $p_{99}$ hat.
Dies ist nur der Fall bei Serialization Failures in \texttt{s2} und \texttt{s3}.
Während es bei \texttt{s2} für erfolgreiche und nicht erfolgreiche Transaktionen der Fall ist, trifft der Fall bei
\texttt{s3} nur bei erfolgreichen Transaktionen zu.

Der Overhead~\ref{fig:s-overhead} verhält sich ähnlich wie die durchschnittliche Ausführungszeit mit kleineren Werten.
Deshalb ist haben Deadlocks den höchsten Overhead, und Serialization failures starten gleich und gehen erst ab
\texttt{Strategie2} deutlicher auseinander.
Dazu verhält sich der Overhead nicht erfolgreicher Transaktion parallel zu den von Lock Timeouts.
Die einzige Besonderheit ist für erfolgreiche Transaktionen bei Lock Timeouts, dass der Overhead erst um 1ms steigt von
\texttt{s0} auf \texttt{s1}, dann stark abfällt zu \texttt{s2} und dann wieder um 20ms steigt für \texttt{s3}.
Das Diagramm~\ref{fig:s-throughput} zeigt den Durchsatz während der einzelnen Experimente.
In diesem Fall ist nur \texttt{Experiment3} zu beachten.
Dadurch ist zu sehen, dass der Durchsatz gleichbleibend ist für Deadlocks und Lock Timeouts.
Bei Serialization failure sinkt der Durchsatz leicht.